% GENERATED FILE. Edit canonical lessons, then run npm run generate:books.
% canonical-source-hash: fnv1a64:ded050ca0fb044a9
% canonical-lessons: RU-C336-dobavlyat, RU-C336-vykladyvat, RU-C336-varyonyy, RU-C336-zharenyy, RU-C336-vynimat, RU-R336-review-335-336, RU-R336-review-final-336

\chapter{Serving up}
\label{ch:ru-serving-up}
\hlchaptermodality{336}

\begin{chapteropening}
\textbf{By the end of this chapter:} \emph{I can say what I added, laid out and took out, and describe food as boiled or fried.}

Complete the last lesson of chapter 336: I can say what I added, laid out and took out, and describe food as boiled or fried.
\end{chapteropening}

\section[\texorpdfstring{\ru{добавлять} (dobavlyát)}{добавлять (dobavlyát)}]{\ru{добавлять} (dobavlyát) — to add}

\label{lesson:RU-C336-dobavlyat}



\begin{quote}
Before the new one: say the Russian for to melt, then the Russian for to turn over.
\end{quote}

\subsection*{You'll want to know: \ru{добавлять}}
\textbf{\ru{добавлять}} (\emph{dobavlyát}) — \textquotedblleft{}to add\textquotedblright{}.

\begin{quote}
\ru{В} \ru{суп} \ru{я} \ru{добавил} \ru{немного} \ru{соли}.
\end{quote}

\begin{quote}
\emph{V sup ya dobávil nemnógo sóli.}
\end{quote}

\begin{quote}
I added a little salt to the soup.
\end{quote}

Imperfective; perfective \emph{\ru{добавить}} (\emph{dobávit}), \emph{\ru{добавлю}}. It takes the accusative, or the genitive for \textquotedblleft{}some\textquotedblright{}: \emph{\ru{добавить} \ru{сахару}}.

\begin{cousinweb}
\emph{\ru{добавлять}} is built on \emph{\ru{до}-} and an old verb \emph{\ru{бавить}} \textquotedblleft{}to make, to keep busy\textquotedblright{}, as in \emph{\ru{забава}} \textquotedblleft{}amusement\textquotedblright{}; the link is usually traced, and EtymDB lists no parent.
\end{cousinweb}

\begin{grammarlens}[title={What you've built}]
One more word for telling, explaining and describing at B1.
\end{grammarlens}

\subsection*{Your turn}
\begin{itemize}
\raggedright
  \item \emph{Say it:} \emph{dobavlyát}
  \item \emph{Say it:} \emph{dobavlyát}, once more
  \item \emph{Say it:} a sentence of your own with \emph{dobavlyát}
  \item \emph{Recall:} say the Russian for to melt, then the Russian for to turn over, then say \emph{dobavlyát} again
\end{itemize}

\begin{tcolorbox}[breakable,colback=teal!4,colframe=teal!35!black,title={Before you move on}]
What does \ru{добавлять} mean? (\textquotedblleft{}To add\textquotedblright{}.) Say it once more.
\end{tcolorbox}
\section[\texorpdfstring{\ru{выкладывать} (vykládyvat)}{выкладывать (vykládyvat)}]{\ru{выкладывать} (vykládyvat) — to lay out, to put out}

\label{lesson:RU-C336-vykladyvat}



\begin{quote}
Before the new one: say the Russian for to turn over, then the Russian for to add.
\end{quote}

\subsection*{You'll want to know: \ru{выкладывать}}
\textbf{\ru{выкладывать}} (\emph{vykládyvat}) — \textquotedblleft{}to lay out, to put out\textquotedblright{}.

\begin{quote}
\ru{Она} \ru{выложила} \ru{печенье} \ru{на} \ru{тарелку}.
\end{quote}

\begin{quote}
\emph{Oná výlozhila pechénye na tarélku.}
\end{quote}

\begin{quote}
She put the biscuits out on a plate.
\end{quote}

Imperfective; perfective \emph{\ru{выложить}} (\emph{výlozhit}). In speech also \textquotedblleft{}to tell everything\textquotedblright{}: \emph{\ru{выложить} \ru{всё}} (\emph{výlozhit vsyo}).

\begin{cousinweb}
\emph{\ru{выкладывать}} is built on \emph{\ru{вы}-} \textquotedblleft{}out\textquotedblright{} and the stem \emph{\ru{клад}-} of \emph{\ru{класть}} \textquotedblleft{}to put\textquotedblright{}, an inherited Slavic verb; EtymDB lists no parent.
\end{cousinweb}

\begin{grammarlens}[title={What you've built}]
One more word for telling, explaining and describing at B1.
\end{grammarlens}

\subsection*{Your turn}
\begin{itemize}
\raggedright
  \item \emph{Say it:} \emph{vykládyvat}
  \item \emph{Say it:} \emph{vykládyvat}, once more
  \item \emph{Say it:} a sentence of your own with \emph{vykládyvat}
  \item \emph{Recall:} say the Russian for to turn over, then the Russian for to add, then say \emph{vykládyvat} again
\end{itemize}

\begin{tcolorbox}[breakable,colback=teal!4,colframe=teal!35!black,title={Before you move on}]
What does \ru{выкладывать} mean? (\textquotedblleft{}To lay out, to put out\textquotedblright{}.) Say it once more.
\end{tcolorbox}
\section[\texorpdfstring{\ru{варёный} (varyónyy)}{варёный (varyónyy)}]{\ru{варёный} (varyónyy) — boiled}

\label{lesson:RU-C336-varyonyy}



\begin{quote}
Before the new one: say the Russian for to add, then the Russian for to lay out.
\end{quote}

\subsection*{You'll want to know: \ru{варёный}}
\textbf{\ru{варёный}} (\emph{varyónyy}) — \textquotedblleft{}boiled\textquotedblright{}.

\begin{quote}
\ru{На} \ru{завтрак} \ru{он} \ru{съел} \ru{варёное} \ru{яйцо}.
\end{quote}

\begin{quote}
\emph{Na závtrak on syel varyónoye yaytsó.}
\end{quote}

\begin{quote}
He had a boiled egg for breakfast.
\end{quote}

\emph{\ru{Варёная} \ru{колбаса}} (\emph{Varyónaya kolbasá}) — a soft boiled sausage.

\begin{cousinweb}
\emph{\ru{варёный}} is the old participle of \emph{\ru{варить}} \textquotedblleft{}to boil\textquotedblright{}, inherited from Proto-Slavic \emph{variti} (EtymDB).
\end{cousinweb}

\begin{grammarlens}[title={What you've built}]
One more word for telling, explaining and describing at B1.
\end{grammarlens}

\subsection*{Your turn}
\begin{itemize}
\raggedright
  \item \emph{Say it:} \emph{varyónyy}
  \item \emph{Say it:} \emph{varyónyy}, once more
  \item \emph{Say it:} a sentence of your own with \emph{varyónyy}
  \item \emph{Recall:} say the Russian for to add, then the Russian for to lay out, then say \emph{varyónyy} again
\end{itemize}

\begin{tcolorbox}[breakable,colback=teal!4,colframe=teal!35!black,title={Before you move on}]
What does \ru{варёный} mean? (\textquotedblleft{}Boiled\textquotedblright{}.) Say it once more.
\end{tcolorbox}
\section[\texorpdfstring{\ru{жареный} (zhárenyy)}{жареный (zhárenyy)}]{\ru{жареный} (zhárenyy) — fried, roast}

\label{lesson:RU-C336-zharenyy}



\begin{quote}
Before the new one: say the Russian for to lay out, then the Russian for boiled.
\end{quote}

\subsection*{You'll want to know: \ru{жареный}}
\textbf{\ru{жареный}} (\emph{zhárenyy}) — \textquotedblleft{}fried, roast\textquotedblright{}.

\begin{quote}
\ru{На} \ru{ужин} \ru{была} \ru{жареная} \ru{рыба}.
\end{quote}

\begin{quote}
\emph{Na úzhin bylá zhárenaya rýba.}
\end{quote}

\begin{quote}
There was fried fish for dinner.
\end{quote}

One \emph{\ru{н}}, as a plain adjective; \emph{\ru{жаренный} \ru{на} \ru{масле}} takes two.

\begin{cousinweb}
\emph{\ru{жареный}} is the old participle of \emph{\ru{жарить}} \textquotedblleft{}to fry\textquotedblright{}, from \emph{\ru{жар}} \textquotedblleft{}heat\textquotedblright{}, the same Slavic root as \emph{\ru{гореть}} \textquotedblleft{}to burn\textquotedblright{}; EtymDB lists no parent.
\end{cousinweb}

\begin{grammarlens}[title={What you've built}]
One more word for telling, explaining and describing at B1.
\end{grammarlens}

\subsection*{Your turn}
\begin{itemize}
\raggedright
  \item \emph{Say it:} \emph{zhárenyy}
  \item \emph{Say it:} \emph{zhárenyy}, once more
  \item \emph{Say it:} a sentence of your own with \emph{zhárenyy}
  \item \emph{Recall:} say the Russian for to lay out, then the Russian for boiled, then say \emph{zhárenyy} again
\end{itemize}

\begin{tcolorbox}[breakable,colback=teal!4,colframe=teal!35!black,title={Before you move on}]
What does \ru{жареный} mean? (\textquotedblleft{}Fried, roast\textquotedblright{}.) Say it once more.
\end{tcolorbox}
\section[\texorpdfstring{\ru{вынимать} (vynimát)}{вынимать (vynimát)}]{\ru{вынимать} (vynimát) — to take out}

\label{lesson:RU-C336-vynimat}



\begin{quote}
Before the new one: say the Russian for boiled, then the Russian for fried.
\end{quote}

\subsection*{You'll want to know: \ru{вынимать}}
\textbf{\ru{вынимать}} (\emph{vynimát}) — \textquotedblleft{}to take out\textquotedblright{}.

\begin{quote}
\ru{Мама} \ru{вынула} \ru{пирог} \ru{из} \ru{духовки}.
\end{quote}

\begin{quote}
\emph{Máma výnula piróg iz dukhóvki.}
\end{quote}

\begin{quote}
Mum took the pie out of the oven.
\end{quote}

Imperfective; perfective \emph{\ru{вынуть}} (\emph{výnut}), \emph{\ru{выну}}.

\begin{cousinweb}
\emph{\ru{вынимать}} is built on \emph{\ru{вы}-} \textquotedblleft{}out\textquotedblright{} and the old Slavic root \emph{-\ru{ним}-}, \emph{-\ru{ять}} \textquotedblleft{}to take\textquotedblright{}, seen in \emph{\ru{взять}} and \emph{\ru{понимать}}; EtymDB traces this root to PIE \emph{h\textsubscript{1}em-}.
\end{cousinweb}

\begin{grammarlens}[title={What you've built}]
One more word for telling, explaining and describing at B1.
\end{grammarlens}

\subsection*{Your turn}
\begin{itemize}
\raggedright
  \item \emph{Say it:} \emph{vynimát}
  \item \emph{Say it:} \emph{vynimát}, once more
  \item \emph{Say it:} a sentence of your own with \emph{vynimát}
  \item \emph{Recall:} say the Russian for boiled, then the Russian for fried, then say \emph{vynimát} again
\end{itemize}

\begin{tcolorbox}[breakable,colback=teal!4,colframe=teal!35!black,title={Before you move on}]
What does \ru{вынимать} mean? (\textquotedblleft{}To take out\textquotedblright{}.) Say it once more.
\end{tcolorbox}
\section[(dialogue)]{Step by step, once more}

\label{lesson:RU-R336-review-335-336}



\begin{quote}
Say the Russian for to cool down and to defrost.
\end{quote}

\subsection*{Your turn}
Say these aloud:

\begin{itemize}
\raggedright
  \item to cool down, to defrost, to chill, to melt, to turn over
  \item to add, to lay out, boiled, fried, to take out
\end{itemize}

\begin{tcolorbox}[breakable,colback=teal!4,colframe=teal!35!black,title={Before you move on}]
To cool down? (\textbf{\ru{остывать}}.) To defrost? (\textbf{\ru{размораживать}}.) To chill? (\textbf{\ru{охлаждать}}.) To melt? (\textbf{\ru{растапливать}}.) To turn over? (\textbf{\ru{переворачивать}}.) To add? (\textbf{\ru{добавлять}}.) To lay out? (\textbf{\ru{выкладывать}}.) Boiled? (\textbf{\ru{варёный}}.) Fried? (\textbf{\ru{жареный}}.) To take out? (\textbf{\ru{вынимать}}.)
\end{tcolorbox}
\section[(dialogue)]{Serving up — once more}

\label{lesson:RU-R336-review-final-336}



\begin{quote}
Say the Russian for to add and to lay out.
\end{quote}

\subsection*{Your turn}
\begin{itemize}
\raggedright
  \item \emph{Say it:} to add, to lay out, boiled, fried, to take out
\end{itemize}

\begin{tcolorbox}[breakable,colback=teal!4,colframe=teal!35!black,title={Before you move on}]
To add? (\textbf{\ru{добавлять}}.) To lay out? (\textbf{\ru{выкладывать}}.) Boiled? (\textbf{\ru{варёный}}.) Fried? (\textbf{\ru{жареный}}.) To take out? (\textbf{\ru{вынимать}}.)
\end{tcolorbox}
